\chapter{Model Formulation and Assumptions}

A linear program optimises a linear
objective over a set described by linear constraints. Before any tableau is
pivoted or any feasible region sketched, the economic story must be
translated into such a model: choose the decision variables, write down what
is being maximised or minimised, and turn every scarce resource, contractual
obligation, and policy rule into a linear inequality or equation. This
translation step is where most exam points are won or lost — a wrong model
solved perfectly is still wrong.

This chapter drills exactly that step: the four components of an LP, the
standard formulation recipe, the sign conventions for inequalities, and the
four assumptions (proportionality, additivity, divisibility, certainty)
under which a linear model is an honest description of the economics.
Solving the models — graphically, and later by the simplex method — begins
in the next chapter; here every answer is a complete model statement, not a
number.

\section{Core Rules}

\begin{rulebox}{The four components of an LP}
Every linear program consists of:
\begin{itemize}[leftmargin=1.6em,itemsep=0.3ex]
  \item \textbf{Decision variables}
        $x_1,\dots,x_n$: the quantities the decision maker controls.
        Always define them in words \emph{with units} (``$x_1=$ number of
        tables produced per week'').
  \item \textbf{Linear objective function}
        $z=c_1x_1+\dots+c_nx_n$, to be maximised (profit, revenue, reach)
        or minimised (cost, risk).
  \item \textbf{Linear constraints}
        $a_{i1}x_1+\dots+a_{in}x_n\ \{\le,=,\ge\}\ b_i$: one per scarce
        resource, requirement, or balance condition.
  \item \textbf{Nonnegativity conditions} $x_j\ge 0$: negative
        production, shipments, or investments are meaningless in almost
        every economic story. State them explicitly — they are constraints
        too.
\end{itemize}
\end{rulebox}

\begin{rulebox}{General and standard form}
The general LP is
\[
  \max\ (\text{or } \min)\quad z=c_1x_1+\dots+c_nx_n
  \quad\text{s.t.}\quad
  \sum_{j=1}^{n} a_{ij}x_j\ \{\le,=,\ge\}\ b_i\ (i=1,\dots,m),
  \quad x_j\ge 0 .
\]
The \emph{standard maximum form} has a maximisation objective and only
$\le$ constraints besides nonnegativity:
\[
  \max\ z=c^{\top}x \quad\text{s.t.}\quad Ax\le b,\ x\ge 0 .
\]
Every LP can be brought to this form by the conversions in the sign-rules
box below; being able to do so mechanically is required for the simplex
chapters.
\end{rulebox}

\begin{rulebox}{Formulation recipe}
\begin{enumerate}[leftmargin=1.8em,itemsep=0.3ex]
  \item \textbf{Identify the decisions.} What does the decision maker
        actually choose? One variable per elementary decision (per product,
        per route, per asset), each with units.
  \item \textbf{Write the objective.} Multiply each variable by its per-unit
        contribution (profit, cost, reach, return) and add up. Decide max
        or min from the story.
  \item \textbf{Write one constraint per restriction.} For each resource:
        (usage per unit)$\,\times\,$(units), summed over variables,
        compared with availability. For each requirement: the same, but
        $\ge$ the target. Check that units match on both sides of every
        constraint.
  \item \textbf{Add sign restrictions,} normally $x_j\ge 0$.
  \item \textbf{State the full model} and re-read the story once more:
        every number in the text must appear in the model, and every
        constraint must be traceable to a sentence.
\end{enumerate}
\end{rulebox}

\begin{rulebox}{Sign conventions and conversions}
\begin{itemize}[leftmargin=1.6em,itemsep=0.3ex]
  \item \textbf{Resource limits} (``at most'', ``available'', ``budget'')
        give $\le$; \textbf{minimum requirements} (``at least'',
        ``contract'', ``must contain'') give $\ge$; \textbf{exact balances}
        (``exactly'', ``all of the budget is invested'') give $=$.
  \item \textbf{Max $\leftrightarrow$ min:}\ \
        $\min z$ has the same optimal points as $\max\,(-z)$, and the
        optimal values satisfy $z^{*}=-(-z)^{*}$.
  \item \textbf{Flipping an inequality:} multiplying both sides by $-1$
        turns $\ge$ into $\le$ (and conversely):
        $a^{\top}x\ge b \iff -a^{\top}x\le -b$.
  \item \textbf{Equalities:} $a^{\top}x=b$ is equivalent to the pair
        $a^{\top}x\le b$ and $a^{\top}x\ge b$.
  \item \textbf{Percentage and ratio conditions} must be cleared of the
        variable-dependent total before they are linear. Example: ``good 1
        is at most $60\%$ of total output'' is
        $x_1\le 0.6(x_1+x_2)$, i.e.\ $0.4x_1-0.6x_2\le 0$, i.e.\
        $2x_1-3x_2\le 0$. Never divide by an expression containing
        variables.
\end{itemize}
\end{rulebox}

\begin{rulebox}{The four LP assumptions}
A linear model is only an honest description of the economics if:
\begin{itemize}[leftmargin=1.6em,itemsep=0.3ex]
  \item \textbf{Proportionality}: each variable's
        contribution to objective and constraints is proportional to its
        level — no quantity discounts, no set-up costs, no economies of
        scale, no overtime premia. Fails: piecewise prices, fixed charges,
        learning effects.
  \item \textbf{Additivity}: total contribution is the
        \emph{sum} of the individual contributions — no interaction between
        variables. Fails: overlapping advertising audiences, synergy or
        crowding effects, joint set-ups.
  \item \textbf{Divisibility}: fractional values of the
        variables are meaningful. Fails: aircraft, workers, machines,
        batches — then integer programming is needed; fractional answers
        are fine for rates, flows, litres, tonnes, and money.
  \item \textbf{Certainty}: all coefficients
        $c_j,a_{ij},b_i$ are known constants. Fails: forecast demand,
        random yields, uncertain prices — then one works with expected
        values plus sensitivity analysis, or stochastic models.
\end{itemize}
\end{rulebox}

\begin{rulebox}{Classic model families}
\begin{center}\small
\begin{tabular}{|l|l|l|}
\hline
\textbf{Family} & \textbf{Decision variables} & \textbf{Typical shape}\\
\hline
Production planning & output of each product & max profit; resources $\le$\\
Diet / blending & amount of each ingredient & min cost; nutrients,
  quality $\ge$\\
Transport & shipment on each route & min cost; supply $\le$,
  demand $\ge$\\
Resource allocation & budget or effort per activity & max return;
  budget $\le$, policy limits\\
Portfolio & money in each asset & max return; budget, risk,
  share limits\\
\hline
\end{tabular}
\end{center}
Recognising the family tells you immediately which direction the
inequalities point and whether to maximise or minimise.
\end{rulebox}

\section{Worked Example}

\begin{worked}{From story to model --- the running example of this volume}
\emph{A workshop makes two products. Each unit of product~1 earns
$40$ euros and each unit of product~2 earns $30$. Product~1 needs $2$ hours
of machine time and $1$ hour of labour; product~2 needs $1$ hour of each.
There are $100$ machine hours and $80$ labour hours available this week.
What should the workshop produce?}

\wstep{Step 1. Name the decision variables, with units.}
\[
  x_{1}=\text{units of product 1 produced this week},
  \qquad
  x_{2}=\text{units of product 2}.
\]
\wtrap{Variables must be quantities the decision-maker actually
\emph{chooses}. ``Machine hours used'' is a consequence, not a decision, and
making it a variable produces a model with redundant equations. Write the
unit --- ``units per week'' --- because it settles the units of every
coefficient that follows.}

\wstep{Step 2. Build the objective from the per-unit contributions.}
\[
  \max\ z=40x_{1}+30x_{2} .
\]
\wtrap{Maximise \emph{contribution}, not revenue, and never mix the two.
Fixed costs may be omitted --- they shift $z$ by a constant and cannot
change the optimal $x$ --- but if they are included they must be included
consistently.}

\wstep{Step 3. One constraint per scarce resource, read down the columns.}
Machine time: product~1 uses $2$ hours per unit, product~2 uses $1$:
\[
  2x_{1}+x_{2}\le100 .
\]
Labour: one hour each:
\[
  x_{1}+x_{2}\le80 .
\]
\wtrap{Coefficients run \emph{down} a column, one per product, and each
constraint must be dimensionally consistent: hours per unit, times units,
gives hours, compared against hours available. Transposing the technology
table is the standard modelling error --- and the units check catches it
immediately, since ``units per hour'' times ``units'' is meaningless.}

\wstep{Step 4. Do not forget the sign restrictions.}
\[
  x_{1}\ge0,\qquad x_{2}\ge0 .
\]
These are constraints, not conventions: without them the model happily
``produces'' negative quantities of the less profitable good to free up
resources, and reports an unbounded profit.

\wstep{Step 5. The complete model.}
\[
  \begin{array}{lrcrcl}
    \max & 40x_{1} & + & 30x_{2} & & \\
    \text{s.t.} & 2x_{1} & + & x_{2} & \le & 100 \quad(\text{machine})\\
                & x_{1} & + & x_{2} & \le & 80 \quad(\text{labour})\\
                & \multicolumn{5}{l}{x_{1},x_{2}\ge0}
  \end{array}
\]
\emph{This model is used in every chapter of this volume.} Chapter~2 solves
it graphically, Chapter~3 puts it in standard form, Chapters~4--5 run the
simplex method on it, Chapter~6 breaks it in four different ways,
Chapter~7 dualises it, and Chapter~8 asks what happens when the numbers
change. Learning it once pays back six times.

\wstep{Step 6. Sanity-check the model before solving it.}
Three quick tests, all worth doing:
\begin{itemize}[leftmargin=1.6em,itemsep=0.2ex]
  \item \textbf{Is anything feasible?} $(0,0)$ satisfies everything, so yes.
  \item \textbf{Is it bounded?} Both constraints cap the total, and both
        objective coefficients are positive, so $z$ cannot run away.
  \item \textbf{Does an obvious answer suggest itself?} Product~1 earns
        more per unit but uses twice the machine time --- per machine hour
        it yields $20$ against product~2's $30$. So the answer will not be
        ``all product~1''.
\end{itemize}
That last observation predicts the solution's character before any
computation, and it is the habit that catches modelling errors.

\wstep{Step 7. The four assumptions being made, and when they fail.}
\begin{center}\footnotesize
\begin{tabular}{ll}
  \hline
  assumption & fails when\\
  \hline
  proportionality & there are set-up costs or quantity discounts\\
  additivity & products interact (shared tooling, joint products)\\
  divisibility & you cannot make $17.14$ aircraft\\
  certainty & demand or yields are random\\
  \hline
\end{tabular}
\end{center}
Each failure has its own remedy --- integer programming, non-linear
programming, stochastic programming --- and each is a different subject.
Knowing which assumption a real problem violates is the practical skill this
chapter teaches.
\end{worked}

\section{Practice Problems}

Formulate each model completely: define the variables with units, state the
objective, justify every constraint, and include the nonnegativity
conditions. Do \emph{not} solve — solving begins in the next chapter.

\begin{problems}
  \item A workshop makes tables and chairs. Each table requires $2$ hours
        of carpentry and $1$ hour of finishing; each chair requires $1$
        hour of carpentry and $2$ hours of finishing. Per week, $8$ hours
        of carpentry and $10$ hours of finishing are available. Profit is
        $5$ per table and $4$ per chair. Formulate the LP that maximises
        weekly profit.
  \item For the LP
        \[
          \max\ z=6x_1+8x_2
          \quad\text{s.t.}\quad
          3x_1+4x_2\le 24,\quad 2x_1+x_2\le 10,\quad x_1\ge 2,\quad
          x_1,x_2\ge 0,
        \]
        identify the four components: decision variables, objective
        function with its coefficients, the constraints (classifying each
        as a resource limit or a minimum requirement) with their
        right-hand sides, and the nonnegativity conditions.
  \item A diet is to be assembled from two foods. One unit of food $1$
        costs $2$ and provides $1$ unit of protein and $3$ units of
        vitamins; one unit of food $2$ costs $3$ and provides $2$ units of
        protein and $1$ unit of vitamins. The diet must supply at least
        $8$ units of protein and at least $9$ units of vitamins. Formulate
        the LP that meets the requirements at minimum cost.
  \item A roaster produces exactly $100$ kg of a coffee blend from bean
        $A$ (cost $4$ per kg, aroma score $6$) and bean $B$ (cost $8$ per
        kg, aroma score $9$). The blend's average aroma score must be at
        least $8$. Formulate the minimum-cost blending LP, and simplify
        the aroma constraint.
  \item A firm ships a product from two warehouses to two shops.
        Warehouse $1$ holds $30$ units, warehouse $2$ holds $20$ units;
        shop $1$ needs $25$ units, shop $2$ needs $25$ units. Unit
        shipping costs:
        \begin{center}
        \begin{tabular}{|l|c|c|}
          \hline
           & Shop 1 & Shop 2\\
          \hline
          Warehouse 1 & 2 & 3\\
          Warehouse 2 & 4 & 1\\
          \hline
        \end{tabular}
        \end{center}
        Formulate the LP that meets all demand at minimum shipping cost,
        and comment on whether the supply and demand constraints may be
        written as equalities.
  \item An agency buys TV spots and radio spots. A TV spot costs $6$ and
        reaches $10$ (thousand people); a radio spot costs $2$ and reaches
        $3$ (thousand people). The budget is $60$. At most $8$ TV spots
        can be booked, and the client insists on at least $5$ radio spots.
        Formulate the LP that maximises total reach.
  \item An investor can place up to $100$ (thousand euros) in an equity
        fund with expected return $8\%$ and a bond fund with expected
        return $4\%$; uninvested money earns nothing. Policy rules: at
        least $30$ must go into bonds, and equity may make up at most
        $60\%$ of the total amount invested. Formulate the LP that
        maximises expected return, writing the percentage rule as a linear
        constraint.
  \item A raw material costs $5$ per unit for the first $100$ units
        ordered and $3$ per unit for every unit beyond $100$. Which LP
        assumption does this pricing scheme violate, and why can the total
        purchasing cost not appear as a single term $c\,x$ in a linear
        objective?
  \item An airline's fleet-planning model returns the answer ``buy
        $x_1=2.5$ aircraft''. Which assumption is violated? Give two
        examples of decision variables for which fractional optimal values
        are perfectly acceptable.
  \item Run separately, a TV campaign reaches $10$ (thousand people) per
        spot and an online campaign $4$ per ad; run together, many viewers
        are counted twice, so total reach is \emph{less} than the sum of
        the separate reaches. Which assumption of the LP model does this
        violate, and where in the model does it show up?
  \item A farmer's planting model uses per-hectare yields and autumn
        selling prices as coefficients — but yields depend on the weather
        and prices on next year's market. Which assumption is violated,
        and what do practitioners typically do about it while still using
        an LP?
  \item Rewrite
        \[
          \min\ z=2x_1-3x_2
          \quad\text{s.t.}\quad
          x_1+x_2\ge 4,\quad x_1-2x_2\le 6,\quad x_1,x_2\ge 0
        \]
        in standard maximum form (maximisation, all functional constraints
        $\le$), and state how the optimal value of the new problem relates
        to $z^{*}$.
  \item A juice producer mixes $x_1$ litres of apple juice and $x_2$
        litres of mango juice. Marketing requires that (a) mango be at
        least $20\%$ of the mix, and (b) the mix contain at least twice as
        much apple as mango. Write both conditions as linear constraints
        of the form $a^{\top}x\le 0$, and check that the two conditions
        are consistent.
  \item A dairy makes yogurt and cheese. One batch of yogurt uses $2$
        litres of milk and $1$ hour of labour; one batch of cheese uses
        $8$ litres of milk and $2$ hours of labour. Per day, $80$ litres
        of milk and $30$ hours of labour are available. A supermarket
        contract obliges the dairy to deliver at least $5$ batches of
        yogurt and at least $2$ batches of cheese daily. Profit is $3$ per
        yogurt batch and $9$ per cheese batch. Formulate the LP, and state
        the direction ($\le$ or $\ge$) of each constraint with a one-line
        justification.
  \item A farmer has $50$ hectares of land and $100$ tonnes of
        fertiliser. One unit of wheat needs $1$ hectare and $2$ tonnes of
        fertiliser and yields profit $4$; one unit of corn needs $1$
        hectare and $1$ tonne of fertiliser and yields profit $3$.
        Formulate the profit-maximising LP.
  \item A bakery sells bread at $3$ per loaf and cake at $5$ per piece. A
        loaf of bread uses $1$ kg of flour and $1$ kg of sugar; a cake
        uses $2$ kg of flour and $1$ kg of sugar. In stock are $12$ kg of
        flour and $8$ kg of sugar. Formulate the LP that maximises
        revenue.
  \item A furniture firm makes desks, shelves, and cabinets with unit
        profits $4$, $3$, and $6$. Per unit, machine time is $1$, $1$, and
        $2$ hours ($40$ hours available); labour is $2$, $1$, and $1$
        hours ($50$ hours available); wood is $3$, $2$, and $4$ square
        metres ($90$ available). The market absorbs at most $10$ cabinets.
        Formulate the LP.
  \item A brewery plans next month's production of lager and ale (in
        hectolitres). Per hectolitre, lager needs $2$ units of malt and
        $1$ unit of hops; ale needs $3$ units of malt and $2$ units of
        hops. Available: $60$ units of malt, $36$ units of hops. The sales
        department \emph{forecasts} that at most $25$ hl of lager and at
        most $15$ hl of ale can be sold. Profit is $5$ per hl of lager and
        $7$ per hl of ale. Brewing is technically only possible in whole
        $5$-hl batches, but you are asked for a linear model. Formulate
        the LP, then name the two assumptions that are questionable in
        this story (with reasons) and one assumption that seems safe.
\end{problems}

\section{Solutions}

\begin{solutions}
  \item \textbf{Variables:} $x_1=$ tables produced per week, $x_2=$ chairs
        produced per week.
        \textbf{Objective:} each table contributes $5$, each chair $4$, so
        maximise $z=5x_1+4x_2$.
        \textbf{Constraints:} carpentry used is $2x_1+1x_2$ hours, at most
        $8$ available; finishing used is $1x_1+2x_2$ hours, at most $10$
        available. Units check: hours per unit $\times$ units $=$ hours on
        both sides. Production cannot be negative.
        \[
          \begin{aligned}
            \max\quad & z=5x_1+4x_2\\
            \text{s.t.}\quad & 2x_1+x_2\le 8 &&\text{(carpentry hours)}\\
            & x_1+2x_2\le 10 &&\text{(finishing hours)}\\
            & x_1,\ x_2\ge 0 .
          \end{aligned}
        \]

  \item \textbf{Decision variables:} $x_1,x_2$ (two of them).
        \textbf{Objective function:} $z=6x_1+8x_2$, a maximisation, with
        coefficients $c_1=6$, $c_2=8$.
        \textbf{Constraints:} $3x_1+4x_2\le 24$ (resource limit, RHS
        $24$); $2x_1+x_2\le 10$ (resource limit, RHS $10$); $x_1\ge 2$
        (minimum requirement, RHS $2$ — e.g.\ a delivery contract for
        good $1$).
        \textbf{Nonnegativity:} $x_1,x_2\ge 0$; note $x_1\ge 2$ makes
        $x_1\ge 0$ redundant, but it is still stated as part of the
        model's sign restrictions.

  \item \textbf{Variables:} $x_1,x_2=$ units of food $1$ and food $2$ in
        the diet.
        \textbf{Objective:} cost is $2$ per unit of food $1$ and $3$ per
        unit of food $2$: minimise $z=2x_1+3x_2$.
        \textbf{Constraints:} protein supplied is $1x_1+2x_2$ units and
        must reach at least $8$; vitamins supplied are $3x_1+1x_2$ units
        and must reach at least $9$. Requirements point the inequalities
        as $\ge$.
        \[
          \begin{aligned}
            \min\quad & z=2x_1+3x_2\\
            \text{s.t.}\quad & x_1+2x_2\ge 8 &&\text{(protein)}\\
            & 3x_1+x_2\ge 9 &&\text{(vitamins)}\\
            & x_1,\ x_2\ge 0 .
          \end{aligned}
        \]
        (Check that the model is not vacuous: $x_1=2$, $x_2=3$ supplies
        protein $2+6=8$ and vitamins $6+3=9$ — both requirements are
        attainable.)

  \item \textbf{Variables:} $x_1,x_2=$ kg of bean $A$ and bean $B$ in the
        blend.
        \textbf{Objective:} minimise cost $z=4x_1+8x_2$.
        \textbf{Constraints:} the blend size is fixed: $x_1+x_2=100$
        (``exactly'' $\Rightarrow$ equality). Average aroma is total aroma
        divided by total kg; clearing the denominator (total $=100$ kg):
        \[
          \frac{6x_1+9x_2}{100}\ge 8
          \iff 6x_1+9x_2\ge 800 .
        \]
        \[
          \begin{aligned}
            \min\quad & z=4x_1+8x_2\\
            \text{s.t.}\quad & x_1+x_2=100 &&\text{(blend size)}\\
            & 6x_1+9x_2\ge 800 &&\text{(average aroma}\ \ge 8)\\
            & x_1,\ x_2\ge 0 .
          \end{aligned}
        \]
        Substituting $x_1=100-x_2$ into the aroma constraint gives
        $600+3x_2\ge 800$, i.e.\ $x_2\ge \frac{200}{3}$: at least two
        thirds of the blend must be bean $B$. The constraint is genuinely
        linear only after the denominator is cleared.

  \item \textbf{Variables:} $x_{ij}=$ units shipped from warehouse $i$ to
        shop $j$, for $i,j\in\{1,2\}$ — four variables, one per route.
        \textbf{Objective:} minimise
        $z=2x_{11}+3x_{12}+4x_{21}+x_{22}$.
        \textbf{Constraints:} each warehouse cannot ship more than it
        holds (supply, $\le$); each shop must receive its demand
        (requirement, $\ge$):
        \[
          \begin{aligned}
            \min\quad & z=2x_{11}+3x_{12}+4x_{21}+x_{22}\\
            \text{s.t.}\quad & x_{11}+x_{12}\le 30
              &&\text{(supply warehouse 1)}\\
            & x_{21}+x_{22}\le 20 &&\text{(supply warehouse 2)}\\
            & x_{11}+x_{21}\ge 25 &&\text{(demand shop 1)}\\
            & x_{12}+x_{22}\ge 25 &&\text{(demand shop 2)}\\
            & x_{11},x_{12},x_{21},x_{22}\ge 0 .
          \end{aligned}
        \]
        Total supply is $30+20=50$ and total demand is $25+25=50$: the
        problem is \emph{balanced}. Meeting all demand then forces every
        warehouse to ship out its whole stock, so all four constraints
        hold with equality at any feasible point and may equivalently be
        written as equalities.

  \item \textbf{Variables:} $x_1=$ number of TV spots, $x_2=$ number of
        radio spots booked.
        \textbf{Objective:} reach (in thousands) is $10$ per TV spot and
        $3$ per radio spot: maximise $z=10x_1+3x_2$.
        \textbf{Constraints:} spending is $6x_1+2x_2$, at most the budget
        $60$; the booking cap gives $x_1\le 8$; the client's insistence
        gives the requirement $x_2\ge 5$.
        \[
          \begin{aligned}
            \max\quad & z=10x_1+3x_2\\
            \text{s.t.}\quad & 6x_1+2x_2\le 60 &&\text{(budget)}\\
            & x_1\le 8 &&\text{(TV cap)}\\
            & x_2\ge 5 &&\text{(radio minimum)}\\
            & x_1,\ x_2\ge 0 .
          \end{aligned}
        \]
        (The model is feasible: $x_1=8$, $x_2=5$ costs $48+10=58\le 60$.
        Strictly, spots come in whole numbers — the divisibility
        assumption; we model them as continuous here.)

  \item \textbf{Variables:} $x_1,x_2=$ thousand euros placed in the
        equity fund and the bond fund.
        \textbf{Objective:} expected return is maximised:
        $z=0.08x_1+0.04x_2$.
        \textbf{Constraints:} the budget gives $x_1+x_2\le 100$ (``up
        to'', so $\le$, not $=$); the bond floor gives $x_2\ge 30$; the
        percentage rule is a ratio condition on the \emph{variable} total
        $x_1+x_2$ and must be expanded, not divided:
        \[
          x_1\le 0.6\,(x_1+x_2)
          \iff 0.4x_1-0.6x_2\le 0
          \iff 2x_1-3x_2\le 0 .
        \]
        \[
          \begin{aligned}
            \max\quad & z=0.08x_1+0.04x_2\\
            \text{s.t.}\quad & x_1+x_2\le 100 &&\text{(budget)}\\
            & x_2\ge 30 &&\text{(bond floor)}\\
            & 2x_1-3x_2\le 0 &&\text{(equity share}\ \le 60\%)\\
            & x_1,\ x_2\ge 0 .
          \end{aligned}
        \]
        Check of the share constraint: at $x_1=60$, $x_2=40$ we get
        $2(60)-3(40)=0$ and the equity share is $60/100=60\%$ — exactly
        the boundary, as it should be.

  \item This violates \textbf{proportionality}. Total purchasing cost is
        \[
          K(x)=
          \begin{cases}
            5x, & 0\le x\le 100,\\
            500+3(x-100), & x>100,
          \end{cases}
        \]
        a piecewise-linear function whose per-unit cost is \emph{not}
        constant: the first unit costs $5$, the $101$st costs $3$. A term
        $c\,x$ in a linear objective assumes every unit contributes the
        same amount $c$, which is false here. Because the discount makes
        $K$ concave, the kink cannot be removed by a simple LP trick
        either (minimising a concave cost rewards jumping past the
        breakpoint); such quantity discounts require piecewise or integer
        programming formulations.

  \item \textbf{Divisibility} is violated: aircraft exist only in whole
        units, so the fractional answer $2.5$ is not implementable, and
        naive rounding can be infeasible or far from optimal — the honest
        model is an integer program. Fractional values are perfectly
        acceptable when the variable measures a divisible quantity, e.g.\
        litres of juice in a blend, tonnes of steel shipped, hours of
        machine time allocated, or thousands of euros invested (also:
        production \emph{rates} interpreted as long-run averages).

  \item \textbf{Additivity} is violated, and it shows up in the
        \emph{objective function}. The LP objective $10x_1+4x_2$ assumes
        total reach is the sum of the two campaigns' separate reaches; with
        overlapping audiences the true reach is
        $10x_1+4x_2-(\text{overlap})$, an interaction term depending on
        both variables jointly. No choice of linear coefficients can
        represent this. (The constraints — money spent simply adds — remain
        additive.) Remedies lie outside pure LP: nonlinear models or
        piecewise-linear approximations of the overlap.

  \item \textbf{Certainty} is violated: the model treats the yields
        $a_{ij}$ and prices $c_j$ as known constants, but they are random
        variables at planting time. Standard practice is to formulate the
        LP with \emph{expected} (or conservatively estimated) values, then
        run a \emph{sensitivity analysis}: how much can a price or yield
        change before the optimal plan changes? (Sensitivity analysis is
        treated systematically in a later chapter.) When randomness is
        severe, stochastic programming replaces the LP.

  \item Two conversions are needed. First, $\min z$ becomes
        $\max\,(-z)$: maximise $w=-2x_1+3x_2$. Second, the $\ge$
        constraint is multiplied by $-1$:
        $x_1+x_2\ge 4 \iff -x_1-x_2\le -4$. The $\le$ constraint and the
        nonnegativity conditions stay as they are:
        \[
          \begin{aligned}
            \max\quad & w=-2x_1+3x_2\\
            \text{s.t.}\quad & -x_1-x_2\le -4\\
            & x_1-2x_2\le 6\\
            & x_1,\ x_2\ge 0 .
          \end{aligned}
        \]
        Both problems have the same optimal points, and the optimal
        values satisfy $z^{*}=-w^{*}$.

  \item (a) ``Mango at least $20\%$ of the mix'':
        \[
          x_2\ge 0.2\,(x_1+x_2)
          \iff 0.2x_1-0.8x_2\le 0
          \iff x_1-4x_2\le 0 .
        \]
        (b) ``At least twice as much apple as mango'':
        \[
          x_1\ge 2x_2 \iff 2x_2-x_1\le 0 .
        \]
        \textbf{Consistency:} together they demand
        $2x_2\le x_1\le 4x_2$, which is satisfiable, e.g.\ $x_1=3$,
        $x_2=1$: apple is three times mango ($\ge$ twice), and the mango
        share is $1/4=25\%\ge 20\%$. In share terms the two rules confine
        mango to between $20\%$ and $\frac{1}{3}\approx 33.3\%$ of the
        mix — a nonempty band.

  \item \textbf{Variables:} $x_1,x_2=$ batches of yogurt and cheese made
        per day.
        \textbf{Objective:} maximise profit $z=3x_1+9x_2$.
        \textbf{Constraints and directions:}
        \begin{itemize}[leftmargin=1.6em,itemsep=0.3ex]
          \item $2x_1+8x_2\le 80$ (milk): resource limit, cannot use more
                milk than is delivered, hence $\le$.
          \item $x_1+2x_2\le 30$ (labour): resource limit, hence $\le$.
          \item $x_1\ge 5$, $x_2\ge 2$ (contract): minimum delivery
                obligations, hence $\ge$.
        \end{itemize}
        \[
          \begin{aligned}
            \max\quad & z=3x_1+9x_2\\
            \text{s.t.}\quad & 2x_1+8x_2\le 80 &&\text{(milk, litres)}\\
            & x_1+2x_2\le 30 &&\text{(labour, hours)}\\
            & x_1\ge 5,\ x_2\ge 2 &&\text{(contract)}\\
            & x_1,\ x_2\ge 0 .
          \end{aligned}
        \]
        The contract point $(5,2)$ uses $10+16=26\le 80$ litres of milk
        and $5+4=9\le 30$ labour hours, so the obligations are compatible
        with the resources and the model is feasible.

  \item \textbf{Variables:} $x_1=$ units of wheat, $x_2=$ units of corn
        planted.
        \textbf{Objective:} maximise profit $z=4x_1+3x_2$.
        \textbf{Constraints:} land used is $1x_1+1x_2$ hectares of the
        $50$ available; fertiliser used is $2x_1+1x_2$ tonnes of the
        $100$ available — both resource limits, hence $\le$.
        \[
          \begin{aligned}
            \max\quad & z=4x_1+3x_2\\
            \text{s.t.}\quad & x_1+x_2\le 50 &&\text{(land, ha)}\\
            & 2x_1+x_2\le 100 &&\text{(fertiliser, t)}\\
            & x_1,\ x_2\ge 0 .
          \end{aligned}
        \]

  \item \textbf{Variables:} $x_1=$ loaves of bread, $x_2=$ cakes
        produced.
        \textbf{Objective:} revenue is $3$ per loaf and $5$ per cake:
        maximise $z=3x_1+5x_2$.
        \textbf{Constraints:} flour used is $1x_1+2x_2$ kg of $12$ kg in
        stock; sugar used is $1x_1+1x_2$ kg of $8$ kg in stock.
        \[
          \begin{aligned}
            \max\quad & z=3x_1+5x_2\\
            \text{s.t.}\quad & x_1+2x_2\le 12 &&\text{(flour, kg)}\\
            & x_1+x_2\le 8 &&\text{(sugar, kg)}\\
            & x_1,\ x_2\ge 0 .
          \end{aligned}
        \]

  \item \textbf{Variables:} $x_1,x_2,x_3=$ desks, shelves, cabinets
        produced.
        \textbf{Objective:} maximise $z=4x_1+3x_2+6x_3$.
        \textbf{Constraints:} one per resource — machine time
        $1x_1+1x_2+2x_3$ of $40$ hours; labour $2x_1+1x_2+1x_3$ of $50$
        hours; wood $3x_1+2x_2+4x_3$ of $90$ square metres — plus the
        market cap on cabinets, $x_3\le 10$.
        \[
          \begin{aligned}
            \max\quad & z=4x_1+3x_2+6x_3\\
            \text{s.t.}\quad & x_1+x_2+2x_3\le 40 &&\text{(machine hours)}\\
            & 2x_1+x_2+x_3\le 50 &&\text{(labour hours)}\\
            & 3x_1+2x_2+4x_3\le 90 &&\text{(wood, m}^2)\\
            & x_3\le 10 &&\text{(market cap)}\\
            & x_1,\ x_2,\ x_3\ge 0 .
          \end{aligned}
        \]
        Every number of the story appears exactly once; each row's units
        are consistent (resource per unit $\times$ units $=$ resource).

  \item \textbf{Variables:} $x_1,x_2=$ hectolitres of lager and ale
        brewed next month.
        \textbf{Objective:} maximise profit $z=5x_1+7x_2$.
        \textbf{Constraints:} malt $2x_1+3x_2\le 60$; hops
        $x_1+2x_2\le 36$; forecast sales ceilings $x_1\le 25$,
        $x_2\le 15$ (producing more than can be sold earns nothing).
        \[
          \begin{aligned}
            \max\quad & z=5x_1+7x_2\\
            \text{s.t.}\quad & 2x_1+3x_2\le 60 &&\text{(malt)}\\
            & x_1+2x_2\le 36 &&\text{(hops)}\\
            & x_1\le 25,\ x_2\le 15 &&\text{(forecast demand)}\\
            & x_1,\ x_2\ge 0 .
          \end{aligned}
        \]
        \textbf{Questionable assumptions:} \emph{divisibility} — brewing
        happens in whole $5$-hl batches, so only multiples of $5$ are
        implementable, and an LP answer such as $x_1=18$ would have to be
        adjusted (integer/batch model); \emph{certainty} — the demand
        ceilings $25$ and $15$ are forecasts, not known constants, so the
        right-hand sides are uncertain data. \textbf{A safe assumption:}
        \emph{additivity} of resource use — the malt consumed by lager
        and by ale simply adds up in the brewhouse; there is no
        interaction between the two products' ingredient use.
\end{solutions}

\section{Additional Problems}

\begin{problems}
  \item A firm makes standard and premium phone cases. A standard case
        uses $2$ units of plastic and $1$ minute of machine time; a
        premium case uses $3$ units of plastic and $2$ minutes of machine
        time. Available per hour: $12$ units of plastic and $8$ minutes of
        machine time. Profits are $3$ (standard) and $5$ (premium).
        Formulate the profit-maximising LP.
  \item A pig feed is mixed from two grains. Per kg, grain $1$ costs $3$
        and provides $2$ units of nutrient $A$ and $1$ unit of nutrient
        $B$; grain $2$ costs $2$ and provides $1$ unit of $A$ and $1$
        unit of $B$. The daily ration must provide at least $10$ units of
        $A$ and at least $8$ units of $B$, and a pig can eat at most $12$
        kg per day. Formulate the minimum-cost feed LP.
  \item Two plants supply two markets. Plant $1$ can produce $40$ units,
        plant $2$ can produce $30$; market $1$ demands $35$, market $2$
        demands $30$. Unit shipping costs: plant $1$ to markets $1,2$:
        $1$ and $4$; plant $2$ to markets $1,2$: $3$ and $2$. Formulate
        the minimum-cost transport LP and state whether the problem is
        balanced.
  \item An endowment of $60$ (thousand euros) must be \emph{fully}
        invested in shares (expected return $9\%$, risk score $8$ per
        thousand) and bonds (expected return $5\%$, risk score $2$ per
        thousand). The portfolio's average risk score per thousand euros
        may be at most $4$. Formulate the return-maximising LP and
        simplify the risk constraint.
  \item Starting production of a product — at any positive level —
        requires cleaning and retooling a machine at a fixed cost of
        $50$, on top of a variable cost of $2$ per unit. Which LP
        assumption does the fixed cost violate, and why?
  \item Rewrite
        \[
          \max\ z=x_1+4x_2
          \quad\text{s.t.}\quad
          x_1+x_2=5,\quad x_1-x_2\ge 1,\quad x_1,x_2\ge 0
        \]
        in standard maximum form (all functional constraints $\le$).
  \item An alloy of at least $100$ kg must be produced from scrap type
        $1$ ($60\%$ copper, cost $3$ per kg) and scrap type $2$ ($30\%$
        copper, cost $1$ per kg). The alloy must contain at least $40\%$
        copper. Formulate the minimum-cost LP, writing the copper
        condition as a linear constraint.
  \item For each mini-story, name the LP assumption that fails:
        (a) a machine's output per hour falls during long production runs
        as it overheats;
        (b) an airline assigns whole aircraft to routes;
        (c) baking two products together in one oven uses less energy
        than baking them separately;
        (d) next quarter's selling prices in the objective are estimates
        from the sales team.
\end{problems}

\section{Additional Solutions}

\begin{solutions}
  \item \textbf{Variables:} $x_1,x_2=$ standard and premium cases made
        per hour. \textbf{Objective:} maximise $z=3x_1+5x_2$.
        \textbf{Constraints:} plastic $2x_1+3x_2\le 12$; machine time
        $x_1+2x_2\le 8$; both are resource limits.
        \[
          \begin{aligned}
            \max\quad & z=3x_1+5x_2\\
            \text{s.t.}\quad & 2x_1+3x_2\le 12 &&\text{(plastic)}\\
            & x_1+2x_2\le 8 &&\text{(machine minutes)}\\
            & x_1,\ x_2\ge 0 .
          \end{aligned}
        \]

  \item \textbf{Variables:} $x_1,x_2=$ kg of grain $1$ and grain $2$ in
        the daily ration. \textbf{Objective:} minimise cost
        $z=3x_1+2x_2$. \textbf{Constraints:} nutrient requirements point
        $\ge$, the stomach capacity points $\le$:
        \[
          \begin{aligned}
            \min\quad & z=3x_1+2x_2\\
            \text{s.t.}\quad & 2x_1+x_2\ge 10 &&\text{(nutrient } A)\\
            & x_1+x_2\ge 8 &&\text{(nutrient } B)\\
            & x_1+x_2\le 12 &&\text{(capacity, kg)}\\
            & x_1,\ x_2\ge 0 .
          \end{aligned}
        \]
        Feasibility check: $x_1=2$, $x_2=6$ gives $A$: $4+6=10\ge 10$,
        $B$: $2+6=8\ge 8$, total $8\le 12$ — the requirements fit inside
        the capacity.

  \item \textbf{Variables:} $x_{ij}=$ units shipped from plant $i$ to
        market $j$. \textbf{Objective:} minimise
        $z=x_{11}+4x_{12}+3x_{21}+2x_{22}$.
        \[
          \begin{aligned}
            \min\quad & z=x_{11}+4x_{12}+3x_{21}+2x_{22}\\
            \text{s.t.}\quad & x_{11}+x_{12}\le 40 &&\text{(capacity plant 1)}\\
            & x_{21}+x_{22}\le 30 &&\text{(capacity plant 2)}\\
            & x_{11}+x_{21}\ge 35 &&\text{(demand market 1)}\\
            & x_{12}+x_{22}\ge 30 &&\text{(demand market 2)}\\
            & x_{ij}\ge 0\ \ (i,j=1,2) .
          \end{aligned}
        \]
        Total capacity $40+30=70$ exceeds total demand $35+30=65$: the
        problem is \emph{unbalanced} (excess supply $5$), so the supply
        constraints must stay $\le$ — some plant will not ship its full
        capacity — while demand is met with $\ge$ (at minimum cost it
        will hold with equality, since extra shipping only adds cost).

  \item \textbf{Variables:} $x_1,x_2=$ thousand euros in shares and
        bonds. \textbf{Objective:} maximise $z=0.09x_1+0.05x_2$.
        \textbf{Constraints:} ``fully invested'' forces the equality
        $x_1+x_2=60$. The average risk score is
        $(8x_1+2x_2)/(x_1+x_2)\le 4$; clearing the (positive) total,
        $8x_1+2x_2\le 4(x_1+x_2)$, i.e.\ $4x_1-2x_2\le 0$, i.e.\
        $2x_1\le x_2$: bonds must be at least twice the equity position.
        Using $x_1+x_2=60$ one may also write it as $8x_1+2x_2\le 240$.
        \[
          \begin{aligned}
            \max\quad & z=0.09x_1+0.05x_2\\
            \text{s.t.}\quad & x_1+x_2=60 &&\text{(full investment)}\\
            & 2x_1-x_2\le 0 &&\text{(average risk}\ \le 4)\\
            & x_1,\ x_2\ge 0 .
          \end{aligned}
        \]
        Check: $x_1=20$, $x_2=40$ has average risk
        $(160+80)/60=4$ — exactly the boundary, matching
        $2(20)-40=0$.

  \item \textbf{Proportionality} fails. The true cost of producing $x>0$
        units is $50+2x$, but producing $x=0$ costs $0$: the cost
        function jumps at zero and is not proportional to the level of
        activity — the first unit effectively costs $52$, each further
        unit $2$. No single coefficient $c$ makes $c\,x$ reproduce this,
        because $c\,x$ is forced through the origin with constant slope.
        Fixed-charge situations require a binary ``switch-on'' variable,
        i.e.\ integer programming.

  \item The objective already maximises. Split the equality into two
        inequalities and flip the $\ge$ constraint:
        $x_1+x_2=5 \iff x_1+x_2\le 5$ and $-x_1-x_2\le -5$;
        $x_1-x_2\ge 1 \iff -x_1+x_2\le -1$. Hence
        \[
          \begin{aligned}
            \max\quad & z=x_1+4x_2\\
            \text{s.t.}\quad & x_1+x_2\le 5\\
            & -x_1-x_2\le -5\\
            & -x_1+x_2\le -1\\
            & x_1,\ x_2\ge 0 .
          \end{aligned}
        \]

  \item \textbf{Variables:} $x_1,x_2=$ kg of scrap $1$ and scrap $2$
        used. \textbf{Objective:} minimise $z=3x_1+x_2$.
        \textbf{Constraints:} output size $x_1+x_2\ge 100$ (``at
        least''); copper content: total copper is $0.6x_1+0.3x_2$ kg, and
        it must be at least $40\%$ of the (variable) total weight:
        \[
          0.6x_1+0.3x_2\ge 0.4\,(x_1+x_2)
          \iff 0.2x_1-0.1x_2\ge 0
          \iff 2x_1-x_2\ge 0 .
        \]
        \[
          \begin{aligned}
            \min\quad & z=3x_1+x_2\\
            \text{s.t.}\quad & x_1+x_2\ge 100 &&\text{(alloy size, kg)}\\
            & 2x_1-x_2\ge 0 &&\text{(copper}\ \ge 40\%)\\
            & x_1,\ x_2\ge 0 .
          \end{aligned}
        \]
        Check: $x_1=40$, $x_2=60$ gives $100$ kg containing
        $24+18=42$ kg of copper, a share of $42\%\ge 40\%$, and indeed
        $2(40)-60=20\ge 0$.

  \item (a) \textbf{Proportionality}: output is no longer proportional
        to running time — the per-hour yield coefficient is not constant.
        (b) \textbf{Divisibility}: aircraft are indivisible units, so
        fractional assignments are meaningless.
        (c) \textbf{Additivity}: the joint energy use is less than the
        sum of the separate uses — an interaction between the two
        activities that a sum of individual terms cannot capture.
        (d) \textbf{Certainty}: the objective coefficients are estimates,
        not known constants; the LP treats them as if they were certain.
\end{solutions}
